\let\stop\empty
\documentclass{exam-zh}
\usepackage{siunitx}
\usepackage{tkz-euclide}
\usepackage{diagbox}
\usepackage{lastpage}
\usepackage{caption}
% 缺失数学符号包，原\pi写成\uppi无定义
\usepackage{amsmath,amssymb}
% tkz-euclide依赖tikz
\usepackage{tikz}
% 修正数学字体冲突，times与xits统一设置
\usepackage{newtxtext}
\usepackage{xits}

\examsetup{
  page/size=a4paper,
  paren/show-paren=true,
  paren/show-answer=true,
  fillin/show-answer=false,
  fillin/no-answer-type=none,
}

\everymath{\displaystyle}

\title{2026 年普通高等学校招生全国统一考试 (新课标 I 卷)}
\subject{数学}

\begin{document}
\secret
\maketitle

\begin{center}
    本试卷共 \pageref{LastPage} 页，19 题。全卷满分 150 分。考试用时 120 分钟。
\end{center}

\begin{notice}
  \item 答卷前，考生务必将自己的姓名、考生号、考场号、座位号填写在答题卡上。
  \item 回答选择题时，选出每小题答案后，用铅笔把答题卡上对应题目的答案标号涂黑，如需改动，用橡皮擦干净后，再选涂其他答案标号。回答非选择题时，将答案写在答题卡上。写在本试卷上无效。
  \item 考试结束后，将本试卷和答题卡一并交回。
\end{notice}

\section{选择题：本题共 8 小题，每小题 5 分，共 40 分。
在每小题给出的四个选项中，只有一项是符合题目要求的。}

% 1.
\begin{question}
  样本数据 $6$，$8$，$4$，$5$，$12$ 的中位数为 \paren
  \begin{choices}
    \item $5$
    \item $6$
    \item $8$
    \item $9$
  \end{choices}
\end{question}

% 2.
\begin{question}
  已知平面向量 $\boldsymbol{a}$，$\boldsymbol{b}$ 不共线，且 $2\boldsymbol{a} + y\boldsymbol{b} = x\boldsymbol{a} - 3\boldsymbol{b}$，则 \paren
  \begin{choices}
    \item $x = 2, y = -3$
    \item $x = -2, y = 3$
    \item $x = 2, y = 3$
    \item $x = -2, y = -3$
  \end{choices}
\end{question}

% 3.
\begin{question}
  已知集合 $A = \left\{\sin\frac{7\pi}{6}, \cos\frac{5\pi}{3}, \tan\frac{5\pi}{4}\right\}$，$B = \left\{-\frac{\sqrt{3}}{2}, -\frac{1}{2}, 1\right\}$，则 $A \cap B =$ \paren
  \begin{choices}
    \item $\left\{-\frac{\sqrt{3}}{2}, -\frac{1}{2}\right\}$
    \item $\left\{-\frac{\sqrt{3}}{2}, 1\right\}$
    \item $\left\{-\frac{1}{2}, 1\right\}$
    \item $\left\{-\frac{\sqrt{3}}{2}, -\frac{1}{2}, 1\right\}$
  \end{choices}
\end{question}

% 4.
\begin{question}
  曲线 $y = 5x + 8\ln x$ 在点 $(1, 5)$ 处的切线方程为 \paren
  \begin{choices}
    \item $y = 3x + 2$
    \item $y = 5x$
    \item $y = 8x - 3$
    \item $y = 13x - 8$
  \end{choices}
\end{question}

% 5.
\begin{question}
  已知抛物线 $C_1: y^2 = 2p_1x \, (p_1 > 0)$ 和 $C_2: x^2 = 2p_2y \, (p_2 > 0)$ 均经过点 $(4, 8)$，则 $C_1$ 的焦点与 $C_2$ 的焦点之间的距离为 \paren
  \begin{choices}
    \item $12$
    \item $4\sqrt{5}$
    \item $6$
    \item $\frac{\sqrt{65}}{2}$
  \end{choices}
\end{question}

% 6.
\begin{question}
  已知函数 $f(x) = \frac{x+2}{\mathrm{e}^x + a}$ 的最大值为 $1$，则 $a =$ \paren
  \begin{choices}
    \item $\frac{1}{2}$
    \item $1$
    \item $\frac{3}{2}$
    \item $2$
  \end{choices}
\end{question}

% 7.
\begin{question}
  一百零八塔位于宁夏回族自治区青铜峡市，以其独特的建筑格局和深厚的历史文化闻名遐迩。该塔群共有 $108$ 座塔，依山势自上而下排成 $12$ 行，将第 $i$ 行中塔的座数记为 $a_i \, (i = 1, 2, \cdots, 12)$，其中 $a_1=1$，$a_2=a_3=3$，$a_4=a_5=5$，且 $a_6, a_7, \cdots, a_{12}$ 是一个首项为 $7$，公差为 $2$ 的等差数列。将 $a_1, a_2, \cdots, a_{12}$ 分为 $6$ 组，每组 $2$ 个数，使得每组的 $2$ 个数之和可构成一个项数为 $6$ 且公差为 $d \, (d > 0)$ 的等差数列，则 $d =$ \paren
  \begin{choices}
    \item $2$
    \item $4$
    \item $6$
    \item $8$
  \end{choices}
\end{question}

% 8.
\begin{question}
  设 $U = \{(x_1, x_2, x_3) \mid x_i \in \{-2, -1, 1, 2\}, i = 1, 2, 3\}$ 为空间中 $64$ 个点构成的集合，点 $P(1, 1, 1)$。记样本空间 $\Omega = \complement_{U} \{P\}$。从 $\Omega$ 中随机取一个点，定义随机变量 $X$ 如下：对 $\Omega$ 中的每个点 $A(x_1, x_2, x_3)$，令 $X(A) = x_1 + x_2 + x_3$，则 $X$ 的数学期望为 \paren
  \begin{choices}
    \item $-\frac{1}{21}$
    \item $-\frac{1}{63}$
    \item $0$
    \item $\frac{1}{7}$
  \end{choices}
\end{question}

\section{多选题：本题共 3 小题，每小题 6 分，共 18 分。
在每小题给出的四个选项中，有多项符合题目要求。
全部选对的得 6 分，部分选对的得部分分，有选错的得 0 分。}

% 9.
\begin{question}
  设 $z = 3 + 2\mathrm{i}$，则 \paren
  \begin{choices}
    \item $\overline{z} = 3 - 2\mathrm{i}$
    \item $|z| = 5$
    \item $z^2 = 5 + 12\mathrm{i}$
    \item $\frac{z+3}{z-\mathrm{i}} \in \mathbf{R}$
  \end{choices}
\end{question}

% 10.
\begin{question}
  在空间中，$A$，$B$ 为两个定点，动点 $C$ 到直线 $AB$ 的距离为 $2$，动点 $D$ 到直线 $AB$ 的距离为 $1$。若二面角 $C-AB-D$ 为 $60^\circ$，则 \paren
  \begin{choices}
    \item $\angle CAD \geqslant 60^\circ$
    \item $CD \geqslant \sqrt{3}$
    \item 当 $AB \perp CD$ 时，$CD \perp \text{平面 } ABD$
    \item 当 $AB \perp \text{平面 } ACD$ 时，$AC \perp AD$
  \end{choices}
\end{question}

% 11.
\begin{question}
  已知圆 $C_1: (x+1)^2 + y^2 = 1$，圆 $C_2: (x-1)^2 + y^2 = 1$，圆 $C_3: x^2 + (y-\sqrt{3})^2 = 1$，直线 $l: y = kx + b$ 与 $C_1, C_2, C_3$ 均有两个交点。记 $l$ 被 $C_1, C_2, C_3$ 截得的弦长分别为 $s_1, s_2, s_3$，则 \paren
  \begin{choices}
    \item $k$ 可以取任意实数
    \item 满足 $s_1 = s_2 = s_3$ 的直线 $l$ 共有 3 条
    \item 满足 $s_1 + s_2 + s_3 = 3$ 的直线 $l$ 多于 3 条
    \item 当 $b = 0$ 时，$s_1 + s_2 + s_3$ 的最大值为 $\frac{2\sqrt{21}}{3}$
  \end{choices}
\end{question}

\section{填空题：本题共 3 小题，每小题 5 分，共 15 分。}

% 12.
\begin{question}
  双曲线 $5x^2 - 6y^2 = 1$ 的离心率为 \fillin．
\end{question}

% 13.
\begin{question}
  已知 $f(x) = 2\sin(ax + \theta) \, (a \in \mathbf{Z}, 0 \leqslant \theta < 2\pi)$ 是偶函数，$f(x)$ 在区间 $\left(0, \frac{\pi}{2}\right)$ 单调递增，则 $\theta =$ \fillin， $f\left(\frac{2\pi}{3}\right) =$ \fillin．
\end{question}

% 14.
\begin{question}
  设实数 $q$ 满足：存在数列 $\{a_n\}$，使得对于任意 $n \in \mathbf{N}^*$，均有 $a_1 + a_2 + \cdots + a_{3n} = n^2 + n$，且 $\{a_n\}$ 中有某些连续 9 项 $a_k, a_{k+1}, \cdots, a_{k+8}$ 是公比为 $q$ 的等比数列，则 $q$ 的最大值为 \fillin．
\end{question}

\section{解答题：本题共 5 小题，共 77 分。解答应写出文字说明、证明过程或演算步骤。}

% 15.
\begin{problem}[points = 13]
  如图，在直三棱柱 $A B C - A_1 B_1 C_1$ 中，$\angle A C B = 90^\circ$，$A C = B C$，$D$，$E$ 分别为 $A B$，$A C_1$ 的中点．
  \begin{enumerate}
    \item[(1)] 证明：$D E \parallel \text{平面 } B C C_1 B_1$；
    \item[(2)] 设 $C C_1 = 2$，直线 $D E$ 与平面 $A C C_1 A_1$ 所成的角为 $45^\circ$，求直线 $D E$ 到平面 $B C C_1 B_1$ 的距离．
  \end{enumerate}
  \begin{flushright}
    \pgfmathsetmacro{\a}{sqrt(2)}
\begin{tikzpicture}[x={(-135:0.5cm)},y={(0:1cm)},z={(90:1cm)},scale=1.5,line join=round]
  \coordinate (B) at (\a,\a,0);
  \node[anchor=150] at (B) {$B$};
  \coordinate (A) at (\a,-\a,0);
  \node[anchor=30] at (A) {$A$};
  \coordinate (C) at (0,0,0);
  \node[anchor=90] at (C) {$C$};

  \coordinate (B1) at (\a,\a,2);
  \node[anchor=-150] at (B1) {$B_1$};
  \coordinate (A1) at (\a,-\a,2);
  \node[anchor=-30] at (A1) {$A_1$};
  \coordinate (C1) at (0,0,2);
  \node[anchor=-90] at (C1) {$C_1$};

  \coordinate (D) at (\a,0,0);
  \node[anchor=90] at (D) {$D$};
  \coordinate (E) at (\a/2,-\a/2,1);
  \node[anchor=-60] at (E) {$E$};

  \draw[thick](A)--(A1)--(B1)--(B)--cycle;
  \draw[thick](A1)--(C1)--(B1);
  \draw[thick,dashed](A)--(C)--(B);
  \draw[thick,dashed](C)--(C1)--(A);
  \draw[thick,dashed](E)--(D);
\end{tikzpicture}
  \end{flushright}
\end{problem}

% 16.
\begin{problem}[points = 15]
  已知在 $\triangle A B C$ 中，$A B = 3$，$B C = 2\sqrt{3}$，$\cos B = \frac{\sqrt{3}}{3}$．
  \begin{enumerate}
    \item[(1)] 求 $\cos A$；
    \item[(2)] 设 $D$，$E$ 两点满足：$D$ 在 $B A$ 的延长线上，$D E \parallel B C$，$A E \perp A C$．若 $D E = \sqrt{6}$，求 $C E$．
  \end{enumerate}
\end{problem}

% 17.
\begin{problem}[points = 15]
  设整数 $N \geqslant 2$，某同学用一个球进行投篮练习，至多投篮 $N$ 次，当且仅当投中 1 次时或 $N$ 次均未投中时，停止练习。设该同学每次投中的概率为 $p \, (0 < p < 1)$，各次投中与否相互独立。记 $X$ 为停止练习时该同学的投篮次数。
  \begin{enumerate}
    \item[(1)] 当 $N = 4$，$p = \frac{1}{3}$ 时，求 $X$ 的分布列；
    \item[(2)] 设 $k$，$m$ 均为自然数。
    \begin{enumerate}
        \item[(i)] 当 $k \leqslant N - 1$ 时，求 $P(X > k)$；
        \item[(ii)] 当 $k + m \leqslant N - 1$ 时，证明：$P(X > k + m \mid X > k) = P(X > m)$．
    \end{enumerate}
  \end{enumerate}
\end{problem}

% 18.
\begin{problem}[points = 17]
  已知椭圆 $C: \frac{x^2}{a^2} + \frac{y^2}{b^2} = 1 \, (a > b > 0)$ 的左焦点为 $F(-1, 0)$，离心率为 $\frac{1}{2}$．
  \begin{enumerate}
    \item[(1)] 求 $C$ 的方程；
    \item[(2)] 设 $O$ 为坐标原点，过 $F$ 且斜率大于 $0$ 的动直线 $l$ 与 $C$ 交于 $P$，$Q$ 两点，其中 $Q$ 在第三象限，直线 $P O$ 与 $C$ 的另一个交点为 $R$．
    \begin{enumerate}
        \item[(i)] 若 $\triangle P Q R$ 的面积是 $\triangle P F O$ 的面积的 3 倍，求 $l$ 的方程；
        \item[(ii)] 求 $\tan \angle P Q R$ 的最小值．
    \end{enumerate}
  \end{enumerate}
\end{problem}

% 19.
\begin{problem}[points = 17]
  已知函数 $f(x)$ 的定义域为 $\mathbf{R}$，且当 $x < 0$ 时，$f(x) = 2^x$．对任意 $x_0 \in \mathbf{R}$，定义集合 $D(x_0) = \{d \in \mathbf{R} \mid f(x_0 + d) > f(x_0)\}$．
  \begin{enumerate}
    \item[(1)] 若当 $x \geqslant 0$ 时，$f(x) = 1 - x$，求 $D(-1)$；
    \item[(2)] 若 $f(x)$ 是奇函数，$f(x_1) \leqslant f(x_2)$ 且 $x_1 x_2 \ne 0$，证明：$D(x_1) \subseteq D(x_2)$；
    \item[(3)] 设 $f(x)$ 满足：① 若 $f(x_1) \leqslant f(x_2)$，则 $D(x_1) \subseteq D(x_2)$；② 当 $0 < x < 1$ 时，$f(x) < f(0)$．
    \begin{enumerate}
        \item[(i)] 证明：$f(0) \geqslant 1$；
        \item[(ii)] 证明：$f(x)$ 在区间 $(0, +\infty)$ 单调递增．
    \end{enumerate}
  \end{enumerate}
\end{problem}

\end{document}